\section*{Vorwort}

Das in dieser Arbeit gewählte generische Maskulinum bezieht sich zugleich auf die männliche, die weibliche und andere Geschlechteridentitäten. Zur besseren Lesbarkeit wird auf die gleichzeitige Verwendung männlicher und weiblicher Sprachformen verzichtet. Alle Geschlechteridentitäten werden ausdrücklich mitgemeint, soweit die Aussagen dies erfordern.

Bei der Erstellung dieser Arbeit wurden digitale Hilfsmittel und KI-basierte Systeme unterstützend eingesetzt. Die Vorgaben der FOM Hochschule für Oekonomie und Management, die Leitlinien der Deutschen Forschungsgemeinschaft zur Sicherung guter wissenschaftlicher Praxis (Kodex 2023) sowie die Datenschutzbestimmungen der Datenschutz-Grundverordnung wurden dabei eingehalten. Personenbezogene oder vertrauliche Daten wurden nicht eingegeben.

Eingesetzt wurde Claude Opus (Anthropic), und zwar zur Verbesserung von Formulierungen und Textstruktur, zur Entwicklung und Bewertung von Titel- und Fragevarianten, zur Strukturierung der Kapitel, zur Unterstützung bei Recherche und Prüfung der Literatur, zur Erstellung der Abbildungen und ihrer Umsetzung in \LaTeX{} sowie zur Prüfung auf sprachliche und formale Korrektheit. Alle inhaltlichen Entscheidungen, die Auswahl und Prüfung der Quellen sowie deren Interpretation stammen vom Verfasser. Übernommen wurden nur Inhalte, die zuvor geprüft und bei Bedarf überarbeitet wurden. Die Nutzung wurde dokumentiert, die Unterlagen können auf Anfrage eingesehen werden.
\\[1cm]
{\myOrt}, September 2026

{\myAutor}
